\section{3. Principal Requirements and Design Drivers}\label{principal-requirements-and-design-drivers}

\subsection{3.1. Mission Purpose and Context}\label{mission-purpose-and-system-context}

EWAS is intended to acquire and assess data on potential ballistic missile launches, identify potential threats and deliver warnings. GSEG is responsible for communications, coordination of two LEO spacecraft, ground fusion, event qualification and delivery of warnings to authorized institutions. It is not a fire-control system.

\subsection{3.2. Principal Driving Requirements}\label{principal-design-requirements}

\begingroup\GSEGTableStyle{3}
\begin{longtable}{|P{0.22}|P{0.29}|P{0.49}|}
\hline\rowcolor{TableHeaderBlue}
\GSEGHeaderText{Area} & \GSEGHeaderText{Identifier /\allowbreak{} B limit} & \GSEGHeaderText{Design consequence}\\ \hline
\endhead
\hline\endlastfoot
Automatic selection of two LEO spacecraft & FUN-001/\allowbreak{}002 & GSN + C\&C, current MCC orbital products\\ \hline
Two observations, assessment and warning & FUN-003/\allowbreak{}004/\allowbreak{}005 & DPS and EWS, four assessment outcomes, without prior operator approval\\ \hline
Coordination and warning time & PERF-001 \textbf{2 s}, PERF-002 \textbf{15 s} & two independently measured critical phases\\ \hline
Single site & ARC-001 & common building and communications risks\\ \hline
Independence of the first warning & ARC-002 & MCC and DAS outside the synchronous first-message path\\ \hline
Resilience and availability & REL-001/\allowbreak{}002/\allowbreak{}003: \textbf{≥99.99\%} & redundancy, readiness status, failure model\\ \hline
Data and security & DATA-001–006; SEC-001–005 & event identity, retention ≥5 years, protection and audit\\ \hline
Interfaces & IF-001–003 & SSEG/\allowbreak{}recipients/\allowbreak{}orbital data, ICDs\\*
\end{longtable}
\endgroup

The system requirement \textbf{>99.9\%} applies to the whole of EWAS, not all GSEG equipment; the more stringent \textbf{≥99.99\%} applies to simultaneous readiness of the two GSEG critical functions described in TRS 3.2.
